\documentclass[11pt]{article}
% Final-manuscript source for the currently agreed model-and-theory scope.
% Numerical experiments and the numerical experiment plan are deferred.
% Original manuscript sources are preserved. Compilation is checked separately.
\usepackage[margin=1in]{geometry}
\usepackage[T1]{fontenc}
\usepackage[utf8]{inputenc}
\usepackage{amsmath,amssymb,amsthm,mathtools}
\usepackage{booktabs,longtable,array}
\usepackage[numbers,sort&compress]{natbib}
\usepackage[hidelinks]{hyperref}
\newtheorem{assumption}{Assumption}[section]
\newtheorem{proposition}{Proposition}[section]
\newtheorem{theorem}{Theorem}[section]
\newtheorem{corollary}{Corollary}[section]
\newtheorem{lemma}{Lemma}[section]
\theoremstyle{remark}
\newtheorem{remark}{Remark}[section]
\DeclareMathOperator*{\argmin}{arg\,min}
\title{Wave-Based Intercity Freight Scheduling with Fleet Circulation}
\author{}
\date{}
\begin{document}
\maketitle

\begin{abstract}
Intercity freight scheduling requires coordinating local collection, line-haul transport, and distribution with the locations and availability of a finite fleet. We develop a time-expanded mixed-integer model in which local service capacity depends on the total stop quantity pooled into a wave and the number of vehicles assigned to it. Precomputed capacity tables represent a continuous routing approximation, while cumulative inventories track vehicles across cities and enforce freight precedence. The model permits direct transportation and independent empty repositioning. Within the local approximation, capacity per vehicle is nondecreasing and pooled capacity is superadditive for a fixed service interval. A rolling-horizon procedure implements the schedule while preserving initiated tasks and inherited inventories. To interpret wave and load decisions, we also study a stationary common-wave benchmark with constant arrivals, full-backlog dispatch, and separable batch-service costs. We derive optimal collection and distribution loads conditional on wave size and establish that repeating a cost-minimizing wave is average-cost optimal among admissible successive complete-wave sizes. The benchmark shares the local routing primitive with the scheduling model, while using a different objective and resource assumptions. A calendar-time argument and a relative-value construction verify its dispatch interpretation.
\end{abstract}

\noindent\textbf{Keywords:} intercity freight; wave-dependent capacity;
continuous routing approximation; mixed-integer scheduling;
vehicle circulation; rolling horizon.

\section{Introduction}
\label{sec:introduction}

Intercity freight service combines dispersed collection and distribution with transportation between cities. Pooling shipments can reduce routing effort per shipment, but orders have release times and delivery deadlines. Vehicles and freight must also be available at the right place and time: a local collection wave creates a subsequent line-haul requirement, and a direct trip changes the fleet available in each city. The operational question is how to coordinate these services when local capacity itself depends on the shipments pooled into a wave.

We study finite-horizon scheduling for a two-city system with fixed fleets and given order time windows. Human-driven vehicles perform collection and distribution, while a separate fleet provides automated hub-to-hub transportation, as in autonomous transfer hub architectures \citep{Lee2024}. The flexible-service model also permits human-driven direct transportation and independent empty repositioning. The objective combines vehicle operating costs and penalties for unserved demand. A wave groups stops assigned to a common service interval, city, and shipment direction; its total quantity differs from the load carried by one vehicle.

Continuous routing approximations provide a way to retain this distinction without selecting individual customer tours. The square-root scaling of \citet{Beardwood1959}, spatial partitioning methods \citep{Karp1977}, and fleet-routing approximations \citep{Daganzo1984} motivate the local service-time model. Closely related, \citet{Banerjee2022} connect single-vehicle zone-area maximization to fleet sizing and service-region partitioning for same-day delivery. Their outputs are tactical design decisions. Here fleet sizes and order windows are inputs, and the outputs are dated services, freight flows, and vehicle movements. The shared methodological element is an aggregate routing approximation.

The paper makes two main contributions. First, it integrates wave-dependent local capacity into a finite-fleet scheduling model. The routing approximation yields capacity tables indexed by service duration and vehicle count, which enter a mixed-integer linear program through binary selectors. Cumulative vehicle inventories and freight precedence connect local waves to intercity services, direct trips, and empty movements. For a fixed local service arc, we show that modeled capacity per vehicle is nondecreasing and pooled capacity is superadditive. These properties explain the capacity mechanism used by the formulation.

Second, it provides a restricted analytical interpretation of wave and vehicle-load choices. In a stationary common-wave benchmark, we derive optimal collection and distribution loads conditional on the wave quantity, retaining both vehicle capacity and wave quantity as load bounds. Under constant arrivals, full-backlog dispatch, and separable batch-service costs, repeating a minimizer of the per-shipment objective is average-cost optimal among admissible successive wave sizes. The benchmark shares the routing primitive with the scheduling model but has continuous loads, a lead-time charge, and no resource competition between batches. It supplies a mechanism interpretation rather than an established relaxation or steady-state limit of the MILP. Complete QVI verification and a comparison on a common integer full-reset threshold domain are retained in the appendices.

Rolling execution supports the main scheduling model by carrying initiated tasks and inherited inventories across computational windows. Capacity construction, scheduling, and execution therefore form the main development of the paper; the stationary analysis is a separate explanatory step. Section~\ref{sec:related} reviews related work. Sections~\ref{sec:capacity} and~\ref{sec:model} present service capacities and finite-horizon scheduling with rolling implementation. Section~\ref{sec:analysis} develops the stationary benchmark, and Section~\ref{sec:discussion} concludes.

\section{Related literature}
\label{sec:related}

\subsection{Continuous approximations and spatial service design}

\citet{Beardwood1959} establish an asymptotic relation between Euclidean tour length, the number of stops, and their spatial distribution. Partitioning algorithms \citep{Karp1977} and analytical fleet-routing models with stop limits \citep{Daganzo1984} connect this geometric scaling to spatial service organization. Our capacity tables use a BHH-inspired approximation with handling and stem travel specified separately. The equal division of wave routing work among vehicles remains a finite-tour modeling approximation.

\citet{Banerjee2022} use continuous approximations to maximize a single vehicle's service-zone area, characterize the associated dispatch policies, and calculate maximum-area functions. Those functions provide fleet-size guidelines and inputs to approximate weighted Voronoi partitioning. Their framework separates depot--zone travel, routing within the zone, setup, and per-order service time. Their term \emph{linehaul} denotes the depot--zone connection; here, line-haul denotes intercity hub-to-hub transport, and local connections are called stem travel. The decision levels also differ: they design single-vehicle zones and fleet sizes, whereas we schedule a given finite fleet across two cities. Our within-wave balanced partitioning assumption does not determine persistent vehicle territories.

\subsection{Consolidation and dispatch decisions}

Transportation and inventory costs jointly shape consolidation decisions \citep{Burns1985}; terminals further change the economics of many-to-many freight flows \citep{Daganzo1987}. \citet{CetinkayaBookbinder2003} analyze quantity-based and time-based shipment consolidation through renewal costs. \citet{Wei2023} compare hybrid, time-based, and quantity-based policies with inbound replenishment and outbound dispatch. These studies motivate explicit separation of batch-service cost and accumulation waiting. Coordinated cycles also arise in integer-ratio inventory policies \citep{Roundy1985}, whose performance guarantees depend on their inventory setting.

Departure timing has long been a decision variable: \citet{Newell1971} studies dispatch under smooth time-varying arrivals, and \citet{Klapp2018} formulate a same-day delivery problem in which a vehicle can dispatch ready orders or wait at available decision epochs. Their dispatch waves are decision opportunities; our local waves identify service intervals whose pooled stops determine capacity. Stochastic routing studies additionally retain random arrival locations and finite vehicle resources \citep{Bertsimas1991,Bertsimas1993}. Those features distinguish full routing policies from a scalar backlog-reset model.

Our stationary benchmark uses a prescribed common-wave class and a routing-based service cost optimized over local vehicle loads. Its contribution is a verified connection between that cost and admissible repeated dispatch decisions. Impulse-control QVIs and verification are established methods \citep{LionsPerthame1986,Belak2017}; the appendices give a self-contained verification for the deterministic separable benchmark. The Poisson comparison concerns the same positive integer full-reset thresholds, without asserting optimality over arbitrary stochastic policies.

\subsection{Multistage freight scheduling and vehicle circulation}

Time-expanded service networks connect departure decisions to resource availability. \citet{Scherr2020} study two-tier logistics with mixed autonomous fleets and develop dynamic discretization discovery for the resulting service network. \citet{He2022} develop an exact algorithm for service network design with hub capacity constraints. Their methods address time discretization and network coupling; our formulation instead uses fixed time nodes and precomputed wave capacities.

\citet{Lee2024} optimize autonomous transfer hub networks combining automated middle-mile transport and human-driven endpoint service. We retain that fleet separation while representing local service through wave capacities and explicitly accounting for human-driven direct and empty movements between cities. Cumulative city inventories make relocation affect subsequent local availability, and cargo precedence connects collection, line-haul, and distribution. Rolling-horizon execution applies this model to revealed orders and preserves earlier commitments. Candidate filtering and independent replay support its implementation, with their scope stated alongside the procedure.

\section{Service system and wave-dependent capacities}
\label{sec:capacity}

We first translate local service times into capacities for dated tasks.
These capacities are the inputs to the finite-fleet scheduling model in
Section~\ref{sec:model}. Figure~\ref{fig:service-architecture} distinguishes
the shipment stages from the spatial pooling within a local wave.

\subsection{Service architecture and operational decisions}

Cities $k\in\{1,2\}$ are connected by a bidirectional corridor. Let
$\mathcal D=\{+,-\}$ be the direction set, with
$o(+)=r(-)=1$ and $o(-)=r(+)=2$. The order set is
$\mathcal L=\bigcup_{d\in\mathcal D}\mathcal L^d$, where each class has
one direction. An order class
$\ell\in\mathcal L^d$ travels from $o(d)$ to $r(d)$ and has quantity
$D_\ell\geq0$, release time $s_\ell$, completion deadline $e_\ell$, and
rejection penalty $p_\ell\geq0$ per unit. Quantities are continuous and
divisible: one unit denotes one pickup-stop and one delivery-stop equivalent.
Thus, the formulation does not assert that a particular indivisible parcel
can be split. We use the time nodes $\mathcal T=\{0,\ldots,T\}$,
where $0\leq s_\ell\leq e_\ell\leq T$ are integers and a period lasts
$\Delta t>0$ minutes. Corridor travel takes $\tau\in\mathbb Z_{>0}$
periods, or $\tau_{\mathrm{phys}}=\tau\Delta t$ minutes.

Initially city $k$ contains $N_k$ human-driven vehicles and $A_k$ automated
vehicles, where $N_k,A_k\in\mathbb Z_{\geq0}$. Write $N=N_1+N_2$ and
$A=A_1+A_2$. The respective capacities are
$M>0$ and $M_A>0$. Human-driven vehicles perform local tours and may also
cross the corridor; automated vehicles operate between hubs. A local tour
starts and ends at its city's hub. A three-leg shipment is collected to
the origin hub, transported between hubs, and distributed from the
destination hub. Direct service keeps the shipment on a human-driven
vehicle across cities. The city-inventory model credits that vehicle to
the other city's available fleet at its modeled trip-completion time. Independent empty
human-driven transfers are also permitted, whether or not there is an order
requiring service.

\begin{table}[tbp]
\centering\small
\caption{Main notation and units for finite-horizon scheduling.}
\label{tab:notation}
\begin{tabular}{@{}p{0.25\linewidth}p{0.69\linewidth}@{}}
\toprule
Symbol & Meaning and unit\\
\midrule
\multicolumn{2}{@{}l}{\textit{Sets and inputs}}\\
$\mathcal L^d,\mathcal T$ & Order classes in direction $d$; discrete time nodes\\
$D_\ell,s_\ell,e_\ell$ & Stop-equivalent demand; release and deadline in periods\\
$N_k,A_k;M,M_A$ & Initial human-driven and automated fleets; load limits in stop equivalents\\
$\Delta t,\tau,\tau_E$ & Minutes per period; loaded and empty intercity travel in periods\\
$a_k,b_k,\rho_k$ & In-zone time coefficients; one-way local stem time in minutes\\
$c_H,c_A,c_D,c_E;p_\ell$ & Currency per vehicle-minute; currency per unserved stop equivalent\\
\midrule
\multicolumn{2}{@{}l}{\textit{Wave and scheduling quantities}}\\
$q,v,n=q/v$ & Total wave stops, vehicle count, and average load per vehicle\\
$x,y,w,e$ & Local, automated, direct, and empty vehicle counts on service arcs\\
$u,\widetilde u$ & Binary selectors for local and direct vehicle counts\\
$g,\widehat g,h,z$ & Local, automated, direct, and unserved stop-equivalent quantities\\
\bottomrule
\end{tabular}
\end{table}

\begin{figure}[tbp]
\centering
\setlength{\unitlength}{1mm}
\begin{picture}(146,65)
\put(0,56){\makebox(146,5)[l]{\small\textbf{Service tasks for one order direction}}}
\put(0,43){\framebox(40,9){\small Collection to hub 1}}
\put(53,43){\framebox(40,9){\small AV hub 1 to hub 2}}
\put(106,43){\framebox(40,9){\small Distribution at city 2}}
\put(40,47.5){\vector(1,0){13}}
\put(93,47.5){\vector(1,0){13}}
\put(0,34){\makebox(146,5)[l]{\small Direct: one human-driven vehicle collects, crosses cities, and delivers.}}
\put(0,26){\makebox(146,5)[l]{\small Direct and empty tasks move vehicles between city inventories at completion.}}
\put(0,17){\makebox(146,5)[l]{\small\textbf{Local wave geometry}}}
\put(0,2){\framebox(57,10){\small Total area $S_k$, total stops $q$}}
\put(57,7){\vector(1,0){20}}
\put(77,2){\framebox(69,10){\small $v$ balanced zones: area $S_k/v$, stops $q/v$}}
\end{picture}
\caption{Service architecture and the balanced local-wave approximation.
The direct description specifies tasks and inventory accounting; it does
not specify unresolved hub-to-service-region connections. Reverse-direction
tasks use the same architecture.}
\label{fig:service-architecture}
\end{figure}

\subsection{Local routing approximation and wave geometry}

For $q\geq0$, define
\begin{equation}
 f_k(q)=a_kq+b_k\sqrt q,\qquad a_k>0,\quad b_k\geq0,
 \qquad
 F_k(\theta)=\left(\frac{\sqrt{b_k^2+4a_k\theta}-b_k}{2a_k}\right)^2,
 \quad \theta\geq0.
 \label{mod:local-function}
\end{equation}
The first term represents handling and the second approximates routing time
under the spatial scaling associated with the BHH limit
\cite{Beardwood1959}. The function $F_k$ is the inverse of $f_k$.
Let $\rho_k\geq0$ denote the average one-way stem time between a hub
and a local service region. Our local wave approximation
assumes that a wave of $v$ vehicles shares a total load $q$ across equal-area
zones, giving per-vehicle duration $2\rho_k+f_k(q)/v$.
This zoning rule is an engineering approximation for finite customer
samples, motivated by the asymptotic routing relation rather than a
finite-sample theorem. For $v>0$, the average vehicle load is $n=q/v$;
the wave quantity $q$ and the vehicle count $v$ are distinct.
Quantities $q,D_\ell,M,M_A$ share the same stop-equivalent unit.
The coefficients $a_k$ and $b_k$ have units of minutes per stop equivalent
and minutes per square root of a stop equivalent, respectively;
$\rho_k$ is in minutes, and $f_k(q)$ is aggregate in-zone service time
in minutes. Period durations are converted to physical time by $\Delta t$.

\paragraph{Why the wave quantity enters the routing term.}
For the balanced zoning approximation, let $S_k$ be the fixed service
area and $\bar\beta_k$ a routing-time coefficient incorporating speed
and the distance metric. Assigning each of $v$ vehicles an area $S_k/v$
and $q/v$ uniformly distributed stops gives the per-vehicle in-zone time
\begin{equation}
 a_k\frac qv+\bar\beta_k\sqrt{\frac{S_k}{v}\frac qv}
 =\frac{a_kq+b_k\sqrt q}{v},\qquad
 b_k=\bar\beta_k\sqrt{S_k}.
 \label{eq:wave-geometry}
\end{equation}
This explains the adopted $f_k(q)/v$ expression. By contrast, $f_k(q/v)$
assigns each vehicle's load to the full service area, producing a routing
term $b_k\sqrt{q/v}$. The distinction concerns spatial pooling, not just
load allocation. The area representation explains an aggregate coefficient;
$b_k$ is an input to the model, not an estimate of area or speed reported
here. The separation of stem, in-zone routing, and stop service follows
the continuous-approximation approach illustrated by
\citet{Banerjee2022}, whose fixed single-vehicle zones serve a
different design problem.

\subsection{Service-arc capacities and pooling properties}
\label{subsec:capacity-analysis}

A local wave is indexed by its start, end, city, and order direction,
$(i,j,k,d)$. Waves that share a start but have different end times are
separate waves; their loads are not pooled to obtain an extra routing
economy. Define the mother arc sets
\begin{align}
 L_k&=\left\lceil\frac{2\rho_k+f_k(M)}{\Delta t}\right\rceil,
 &\mathcal A_k&=\{(i,j)\in\mathcal T^2:1\leq j-i\leq L_k\},
 \label{mod:local-arcs}\\
 \mathcal A_A&=\{(i,j)\in\mathcal T^2:j-i=\tau\},
 &\mathcal A_E&=\{(i,j)\in\mathcal T^2:j-i=\tau_E\},
 \label{mod:movement-arcs}\\
 L_D&=\tau+\left\lceil
 \frac{2\max_k\rho_k+f_1(M)+f_2(M)}{\Delta t}\right\rceil,
 &\mathcal A_D&=\{(i,j)\in\mathcal T^2:\tau\leq j-i\leq L_D\}.
 \label{mod:direct-arcs}
\end{align}
Here $\tau_E\in\mathbb Z_{>0}$ is the empty-transfer duration; its default
is $\tau$. These finite domains are part of the dispatch model. In
particular, all tasks finish within the actual operating horizon $T$.

For $v=1,\ldots,N$, precompute the local capacity
\begin{equation}
 K^k_{ij}(v)=\min\{vM,F_k(v[(j-i)\Delta t-2\rho_k]^+)\},
 \qquad [x]^+=\max\{x,0\}.
 \label{mod:local-capacity}
\end{equation}
Fix a local arc and suppress its city index. Write $D=(j-i)\Delta t$,
$B=(D-2\rho)^+$, and $K(v)=\min\{vM,f^{-1}(vB)\}$, extended to
$v>0$, with $K(0)=0$. The load constraint $q/v\le M$ and the duration
constraint $2\rho+f(q)/v\le D$ define the same capacity for positive $q$.

\begin{proposition}[Capacity economies within a fixed service arc]
\label{prop:analysis-capacity}
For fixed $D,\rho,M$, the ratio $K(v)/v$ is nondecreasing for $v>0$.
Consequently,
\begin{equation}
 K(v_1+v_2)\ge K(v_1)+K(v_2),\qquad v_1,v_2\ge0.
 \label{eq:analysis-superadditive}
\end{equation}
These statements apply in particular to integer fleet sizes. They do not
imply that $K$ is convex or that either inequality is strict.
\end{proposition}

Pooling within the same city, direction, start, and end time weakly
increases modeled capacity relative to separate waves of equal duration.
The gain can disappear when load limits bind. The derivation and proof
are in Appendix~\ref{app:proofs}.

\subsection{Direct-service capacity and route conventions}

For a direct wave in direction $d$, the origin collection is zoned, whereas
each vehicle's destination load is spread over the destination service
area. We adopt the following aggregate direct-service duration approximation:
\begin{equation}
 \Psi_d(q,v)=2\rho_{o(d)}+
 \frac{f_{o(d)}(q)}{v}+\tau\Delta t+f_{r(d)}(q/v).
 \label{mod:direct-duration}
\end{equation}
Set $\widetilde K^d_{ij}(v)$ equal to the largest $q\in[0,vM]$ satisfying
$\Psi_d(q,v)\leq(j-i)\Delta t$, or zero if this set is empty. We retain
the convention $\widetilde K^d_{ij}(v)=0$ when $j-i=\tau$.
The table also has a closed form. For
\begin{equation}
 B_{ij}^d=(j-i-\tau)\Delta t-2\rho_{o(d)},\qquad
 \alpha_{dv}=\frac{a_{o(d)}+a_{r(d)}}{v},\qquad
 \beta_{dv}=\frac{b_{o(d)}}{v}+\frac{b_{r(d)}}{\sqrt v},
 \label{mod:direct-coefficients}
\end{equation}
if $j-i>\tau$ and $B_{ij}^d>0$, then
\begin{equation}
 \widetilde K^d_{ij}(v)=\min\left\{vM,
 \left(\frac{\sqrt{\beta_{dv}^2+4\alpha_{dv}B_{ij}^d}
                   -\beta_{dv}}{2\alpha_{dv}}\right)^2\right\};
 \label{mod:direct-capacity}
\end{equation}
otherwise the capacity is zero. Unlike a symmetric pooled-tour formula,
\eqref{mod:direct-duration} generally gives different capacities in the two
directions.



\section{Finite-horizon scheduling and rolling implementation}
\label{sec:model}

The capacity tables from Section~\ref{sec:capacity} determine a time-expanded
MILP for fixed initial fleets and given order windows. The same formulation
is then used with revealed orders and inherited commitments in rolling execution.

\subsection{Decisions and objective}

For each applicable arc and direction, let $x_{ij}^{kd}$ count the human-driven
vehicles assigned to a local wave, $y_{ij}^d$ the automated vehicles,
$w_{ij}^d$ the human-driven vehicles assigned to a direct wave, and
$e_{ij}^d$ the independent empty transfers. All four are nonnegative integers. Binary selectors
$u_{ijv}^{kd}$ and $\widetilde u_{ijv}^d$ identify a positive local or
direct vehicle count $v$. Zero vehicles are represented by selecting no positive
size. For $\ell\in\mathcal L^d$, the nonnegative continuous variables
$g_{ij}^{k\ell}$, $\widehat g_{ij}^{\ell}$, and $h_{ij}^{\ell}$ are the
local, automated, and direct quantities of that order. The unserved quantity
is $0\leq z_\ell\leq D_\ell$.

Let $c_H,c_A,c_D,c_E\geq0$ be costs in currency per vehicle-minute.
When hourly rates are supplied to the implementation, they are divided by
$60$ once. With zero hub-transfer delay and no separate transfer fee, the
objective is
\begin{align}
 \min\quad &\sum_{\ell\in\mathcal L}p_\ell z_\ell
 +c_H\Delta t\sum_{k,d}\sum_{(i,j)\in\mathcal A_k}(j-i)x_{ij}^{kd}
 +c_A\Delta t\sum_d\sum_{(i,j)\in\mathcal A_A}(j-i)y_{ij}^d
 \nonumber\\
 &+c_D\Delta t\sum_d\sum_{(i,j)\in\mathcal A_D}(j-i)w_{ij}^d
 +c_E\Delta t\sum_d\sum_{(i,j)\in\mathcal A_E}(j-i)e_{ij}^d.
 \label{mod:objective}
\end{align}
The default empty-transfer cost is $c_E=c_D$, although it can be specified
separately. The unit of $p_\ell$ is currency per stop equivalent, so
both rejection and vehicle-time terms are monetary costs. The objective
prices vehicle operation and rejection; accumulation waiting, shipment
lead time, and sorting delay are outside this finite-horizon objective.
Vehicle cost is charged for the entire assigned task duration,
including any slack within that interval. Unserved-demand penalties allow
rejection, rather than imposing compulsory acceptance of every order.

\subsection{Wave and line-haul capacity}

The local constraints, for every $k,d,(i,j)\in\mathcal A_k$, are
\begin{equation}
 \begin{aligned}
 x_{ij}^{kd}&=\sum_{v=1}^N v u_{ijv}^{kd},\qquad
 \sum_{v=1}^N u_{ijv}^{kd}\leq1,\\
 \sum_{\ell\in\mathcal L^d}g_{ij}^{k\ell}
 &\leq\sum_{v=1}^N K^k_{ij}(v)u_{ijv}^{kd}.
 \end{aligned}
 \label{mod:local-links}
\end{equation}
For each $d,(i,j)\in\mathcal A_D$, impose
\begin{equation}
 \begin{aligned}
 w_{ij}^d&=\sum_{v=1}^N v\widetilde u_{ijv}^d,\qquad
 \sum_{v=1}^N\widetilde u_{ijv}^d\leq1,\\
 \sum_{\ell\in\mathcal L^d}h_{ij}^{\ell}
 &\leq\sum_{v=1}^N\widetilde K^d_{ij}(v)\widetilde u_{ijv}^d.
 \end{aligned}
 \label{mod:direct-links}
\end{equation}
Automated capacity is shared across all orders on an arc:
\begin{equation}
 \sum_{\ell\in\mathcal L^d}\widehat g_{ij}^{\ell}
 \leq M_A y_{ij}^d,
 \qquad d\in\mathcal D,\ (i,j)\in\mathcal A_A.
 \label{mod:auto-capacity}
\end{equation}
Both selector ranges use the total human-driven fleet $N$, not just its
initial city allocation $N_k$: a vehicle can migrate and subsequently join
a local wave. The inventory constraints below enforce the actual
availability in each city. Neither $e$ nor $y$ requires a positive customer
load. Nor does \eqref{mod:direct-links} force $w$ to carry cargo: a
$\tau$-period zero-cargo direct movement remains admissible under the above
convention. The separate variable $e$ makes empty relocation available
independently of customer candidates and permits a distinct cost or duration.

\subsection{Vehicle conservation in time and space}

Let $d_k$ be the direction leaving city $k$ and $\bar d_k$ its opposite.
For every $k$ and $t\in\mathcal T$, define the idle human-driven inventory
\begin{align}
 I^H_k(t)={}&N_k
 -\sum_{(i,j)\in\mathcal A_D:i\leq t}w_{ij}^{d_k}
 -\sum_{(i,j)\in\mathcal A_E:i\leq t}e_{ij}^{d_k}
 +\sum_{(i,j)\in\mathcal A_D:j\leq t}w_{ij}^{\bar d_k}
 +\sum_{(i,j)\in\mathcal A_E:j\leq t}e_{ij}^{\bar d_k}
 \nonumber\\
 &-\sum_{d\in\mathcal D}\sum_{(i,j)\in\mathcal A_k:i\leq t<j}x_{ij}^{kd}
 \ \geq0.
 \label{mod:human-inventory}
\end{align}
The corresponding automated inventory is
\begin{equation}
 I^A_k(t)=A_k-
 \sum_{(i,j)\in\mathcal A_A:i\leq t}y_{ij}^{d_k}
 +\sum_{(i,j)\in\mathcal A_A:j\leq t}y_{ij}^{\bar d_k}
 \geq0.
 \label{mod:auto-inventory}
\end{equation}
Local occupancy is half-open, $i\leq t<j$, so a vehicle completing at $j$
can start a new task at that same node. The same convention allows an
intercity arrival and a departure at the same node. Summing the two
automated inventory inequalities also gives the implied fleet bound
\begin{equation}
 \sum_d\sum_{(i,j)\in\mathcal A_A:i\leq t<j}y_{ij}^d\leq A.
 \label{mod:auto-busy}
\end{equation}
For human-driven vehicles, \eqref{mod:human-inventory} simultaneously
accounts for local occupancy, migration, and independent empty transfers;
a fixed per-city busy-vehicle bound alone would not do so.

\subsection{Order windows, precedence, and demand balance}

For $\ell\in\mathcal L^d$, impose the forbidden combinations
\begin{align}
 g_{ij}^{o(d)\ell}&=0 &&\text{if }i<s_\ell,
 &g_{ij}^{r(d)\ell}&=0 &&\text{if }j>e_\ell,
 \label{mod:local-windows}\\
 h_{ij}^{\ell}&=0 &&\text{if }i<s_\ell\text{ or }j>e_\ell.
 \label{mod:direct-windows}
\end{align}
Delivery is credited at the end of the destination tour or direct trip.
For each $t\in\mathcal T$, cumulative precedence requires
\begin{align}
 \sum_{(i,j)\in\mathcal A_A:i\leq t}\widehat g_{ij}^{\ell}
 &\leq\sum_{(i,j)\in\mathcal A_{o(d)}:j\leq t}g_{ij}^{o(d)\ell},
 \label{mod:pickup-before-auto}\\
 \sum_{(i,j)\in\mathcal A_{r(d)}:i\leq t}g_{ij}^{r(d)\ell}
 &\leq\sum_{(i,j)\in\mathcal A_A:j\leq t}\widehat g_{ij}^{\ell}.
 \label{mod:auto-before-delivery}
\end{align}
Goods can wait at either hub, but cannot depart before arrival. Equality
of event times permits immediate handover; no sorting queue is modeled in
this finite-horizon formulation. Finally, require
\begin{equation}
 \sum_{(i,j)\in\mathcal A_k}g_{ij}^{k\ell}
 +\sum_{(i,j)\in\mathcal A_D}h_{ij}^{\ell}+z_\ell=D_\ell,
 \qquad k\in\{o(d),r(d)\},\quad \ell\in\mathcal L^d.
 \label{mod:demand}
\end{equation}
At $t=T$, \eqref{mod:pickup-before-auto}--\eqref{mod:demand} imply that
the total transshipment quantity is identical at collection, automated
line-haul, and distribution. Consequently, an order cannot receive rejection
credit for goods that have already been collected and left unfinished.
The direct and transshipment quantities are separate portions of the same
demand and are not counted twice.

Equations \eqref{mod:objective}--\eqref{mod:demand}, the arc domains, binary
selectors, nonnegative integer vehicle counts, and nonnegative continuous
goods and rejection variables define a mixed-integer linear program once
the capacity tables have been computed. Transfer-only operation with city-dedicated human-driven fleets is the
restriction $w=e=h=0$; the direct-wave selectors are then zero as well.
This restriction closes both direct freight service and human-driven empty
intercity repositioning. Comparing it with the flexible model therefore
changes both freight routing and human-driven fleet mobility.
The core formulation places no prescribed lower or upper share on direct
service. Positive transfer delays, transfer fees, and direct-share bounds
are optional implementation extensions and are not used in the model
analyzed here.

\subsection{Service variants and the physical endpoint}

The default model leaves the final city allocation free. This is a distinct
choice from allowing unfinished tasks: every modeled arc still ends by
$T$. If an actual shutdown requires human-driven terminal stocks
$\bar N_1+\bar N_2=N$, it can instead impose
\begin{equation}
 I^H_k(T)=\bar N_k,\qquad k=1,2.
 \label{mod:optional-terminal}
\end{equation}
Such a constraint must account for both $w$ and $e$. A computational
rolling-window boundary is not an operating shutdown and does not justify
\eqref{mod:optional-terminal}. The rolling implementation described in
Section~\ref{sec:method}
uses the free terminal allocation; enforcing a genuine future terminal
target would additionally require a feasible completion policy for all
previous commitments.

\subsection{Rolling-horizon implementation}
\label{sec:method}

The offline model uses all given orders. Rolling execution instead uses
revealed orders and commits near-term actions. Its inputs are the model
data, a prediction length $H\ge1$, a control step $h\in\{1,\ldots,H\}$,
a completion extension $E\ge0$ (all in periods), and a total run budget.
At decision node $t$, define
\begin{equation}
 b_t=\min\{T,t+H\},\qquad c_t=\min\{T,b_t+E\},\qquad
 a_t=\min\{T,t+h\}.
 \label{alg:window}
\end{equation}
The revealed set $\mathcal L(t)=\{\ell:s_\ell\le t\}$ retains orders after
their deadlines so that collected goods and vehicle commitments remain
accounted for. All times are absolute.

Every task started before $t$ is fixed in full, including vehicle counts,
selectors, and cargo. Missing historical starts are fixed to zero.
Completed trips remain in cumulative balances. New tasks satisfy
\begin{align}
 t\leq i<b_t,\quad j\leq c_t
 &\qquad\text{for origin collection, direct trips, and empty transfers},
 \label{alg:origin-cut}\\
 t\leq i<c_t,\quad j\leq c_t
 &\qquad\text{for automated line-haul and destination distribution}.
 \label{alg:downstream-cut}
\end{align}
The extension allows downstream completion beyond the new-origin admission
window. These restrictions apply after history has been fixed: ongoing
tasks are not truncated at $b_t$ or $c_t$, and every task respects the
physical endpoint $T$ (Figure~\ref{fig:rolling-timeline}).

\begin{enumerate}
 \item Reveal orders and reconstruct absolute vehicle and cargo history;
       optionally filter new cargo candidates.
 \item Build the MILP for known orders, fix historical starts and zeros,
       and apply \eqref{alg:origin-cut}--\eqref{alg:downstream-cut} to new
       tasks. Solve within the remaining total budget.
 \item With an incumbent, commit all vehicle, selector, and cargo decisions
       starting in $[t,a_t)$, retaining whole durations and every positive
       cargo quantity. Advance to $a_t$.
 \item With no incumbent or an exhausted budget, report partial execution.
       At $t=T$, reconstruct final rejection from completed deliveries.
\end{enumerate}

Each incumbent must complete the goods it collects, as required by
\eqref{mod:demand}; provisional rejection is not a committed physical
event. Computational window endpoints impose no city-return condition.
The procedure preserves past commitments but has no general
recursive-feasibility guarantee: if a continuation cannot be found, its
partial status remains visible.

Appendix~\ref{app:implementation} gives the cargo filter and independent
replay. Filtering retains vehicle domains and historical cargo, while replay
checks original service-time inequalities, inventories, and goods balances.
Costs count each initiated task once. Window bounds and gaps apply to
individual subproblems, not to the full rolling execution; partial execution
has no certified complete-horizon objective.

\begin{figure}[tbp]
\centering
\setlength{\unitlength}{1mm}
\begin{picture}(146,44)
\put(0,10){\vector(1,0){146}}
\multiput(17,9)(0,2){15}{\line(0,1){1}}
\multiput(44,9)(0,2){15}{\line(0,1){1}}
\multiput(105,9)(0,2){15}{\line(0,1){1}}
\put(140,8){\line(0,1){31}}
\put(5,4){\makebox(0,4){$i$}}
\put(17,4){\makebox(0,4){$t$}}
\put(44,4){\makebox(0,4){$b_t$}}
\put(74,4){\makebox(0,4){$j$}}
\put(105,4){\makebox(0,4){$c_t$}}
\put(140,4){\makebox(0,4){$T$}}
\put(5,25){\line(1,0){69}}
\put(5,23){\line(0,1){4}}
\put(74,23){\line(0,1){4}}
\put(5,28){\makebox(69,5){\small Inherited task, fixed in full}}
\put(76,17){\makebox(29,6){\small Vehicle available}}
\put(5,0){\makebox(136,3){\scriptsize Schematic absolute times: $i<t<b_t<j\le c_t\le T$}}
\end{picture}
\caption{An inherited task continues beyond the current origin-admission
window $b_t$ until its fixed completion time $j$. The real operating
endpoint is $T$. This is a mechanism diagram, not an experimental schedule.}
\label{fig:rolling-timeline}
\end{figure}

\section{A stationary benchmark for wave and load decisions}
\label{sec:analysis}

This auxiliary benchmark uses the local routing mechanism to interpret wave
and load choices. Table~\ref{tab:model-scope} records its differences from the
scheduling problem. It is not an established relaxation, equivalent
reformulation, or steady-state limit of the MILP.

\begin{table}[tbp]
\centering\small
\caption{The scheduling model and the auxiliary stationary benchmark.}
\label{tab:model-scope}
\begin{tabular}{@{}p{0.20\linewidth}p{0.37\linewidth}p{0.35\linewidth}@{}}
\toprule
Dimension & Scheduling model & Stationary benchmark\\
\midrule
Decisions & Dated arcs, vehicle counts, cargo flows & Wave quantity and continuous vehicle loads\\
Information & Given orders offline; revealed orders when rolling & Constant fluid arrivals per direction\\
Resources & Fixed integer fleets and city inventories & No competition between batches for fleets or sorting\\
Service & Transfer, direct, and independent empty tasks & Three-leg common-wave policy class\\
Objective & Vehicle costs and unserved-demand penalties & Per-shipment cost including mean lead time\\
Delivery & End of the destination service arc & Mean delivery during a distribution tour\\
\bottomrule
\end{tabular}
\end{table}

\subsection{Benchmark assumptions and cost accounting}
\label{subsec:benchmark}

We use the following benchmark to isolate the dispatch decision. All physical
times in this section are in minutes; $\tau_{\mathrm{phys}}=\tau\Delta t$ is
the highway travel time, rather than a number of time periods.

\begin{assumption}[Stationary common-wave benchmark]
\label{ass:analysis-benchmark}
The two cities are identical. Shipments in each direction arrive as a
deterministic fluid at rate $\Lambda>0$ shipments per minute, with uniformly
distributed origins and destinations. Every shipment is served through the
three-leg chain. The three legs share a common wave of size $q>0$, with
headway $q/\Lambda$. Vehicle loads are continuous, and fleets are sized to
the offered workload, without vehicle queues or inter-wave resource
competition. Sorting has unlimited capacity, collection tours have no
additional slowest-vehicle delay, and no hard delivery deadline is imposed.
Common-wave operation is the prespecified policy class under study,
rather than a conclusion that independently controlled stages synchronize
optimally.
\end{assumption}

Let $n_1$ and $n_2$ denote the collection and distribution loads, respectively.
The feasible load domain is
\begin{equation}
 0<n_1,n_2\le\min\{M,q\},\qquad
 s(q)=a+\frac b{\sqrt q}.
 \label{eq:analysis-load-domain}
\end{equation}
Here $a>0$, $b>0$, $\rho>0$, and $M_A\ge M>0$. We take $c_H>0$ and
$c_A\ge0$ as the manual and automated vehicle costs per minute, and
$\omega>0$ as the cost per shipment-minute of lead time. Each hub transfer
has fixed time $\theta\ge0$ and fee $h_{\mathrm{hand}}\ge0$ per shipment.
As in the steady-state approximation, the highway cost per shipment is
$c_A\tau_{\mathrm{phys}}/\min\{q,M_A\}$. For waves exceeding $M_A$, this
treats line-haul vehicle requirements continuously and omits integer
vehicle-rounding costs.

The generalized cost per shipment in one direction is
\begin{align}
 G(q;n_1,n_2)={}&c_H\left[2s(q)+\frac{2\rho}{n_1}
                                      +\frac{2\rho}{n_2}\right]
 +\frac{c_A\tau_{\mathrm{phys}}}{\min\{q,M_A\}}
 \notag\\
 &+\omega\left[\frac q{2\Lambda}
       +s(q)\left(n_1+\frac{n_2}2\right)
       +3\rho+\tau_{\mathrm{phys}}+2\theta\right]
 +2h_{\mathrm{hand}}.
 \label{eq:analysis-generalized-cost}
\end{align}
The term $q/(2\Lambda)$ is the mean wait for a wave to form. All collected
shipments remain in the collection chain until its tour returns, whereas a
shipment on a distribution tour is delivered after the outward stem and half
the in-zone tour on average. This explains the distinct coefficients of
$n_1$ and $n_2$. In particular, this benchmark measures delivery during a
distribution tour. Section~\ref{sec:model} instead registers delivery at the
end of a service arc. The benchmark also adds a lead-time cost that is absent
from the base finite-horizon vehicle-cost and rejection-penalty objective.

\subsection{Optimal vehicle loads conditional on a wave}

\begin{proposition}[Optimal loads conditional on a wave]
\label{prop:analysis-loads}
Under Assumption~\ref{ass:analysis-benchmark}, the minimizing loads for each
$q>0$ are unique and are
\begin{align}
 n_1^\circ(q)&=\min\left\{M,q,
                   \sqrt{\frac{2c_H\rho}{\omega s(q)}}\right\},
 \label{eq:analysis-n1}\\
 n_2^\circ(q)&=\min\left\{M,q,
                   \sqrt{\frac{4c_H\rho}{\omega s(q)}}\right\}.
 \label{eq:analysis-n2}
\end{align}
Consequently $n_1^\circ(q)\le n_2^\circ(q)\le\sqrt2\,n_1^\circ(q)$.
\end{proposition}

The clipping by $q$ in these formulas is necessary: a vehicle cannot carry
more shipments than the entire wave. The load asymmetry reflects the different
lead-time exposure of collection and distribution, rather than a difference
in their physical vehicle capacities. Define
\begin{equation}
 G^\circ(q)=G(q;n_1^\circ(q),n_2^\circ(q)),\qquad
 C(q)=q\left[G^\circ(q)-\frac{\omega q}{2\Lambda}\right].
 \label{eq:analysis-dispatch-cost}
\end{equation}
The function $C$ is the generalized service cost of a dispatched batch. It
includes post-dispatch lead time, but excludes the waiting already incurred
while the batch forms.

\subsection{Optimal repeated dispatch under separable batch costs}
\label{subsec:dynamic-wave}

The dispatch model permits successive common waves to have different sizes,
while retaining the same cost and physical approximations.

\begin{assumption}[Separable service accounting]
\label{ass:analysis-accounting}
A dispatch of size $q$ incurs the prepaid service cost $C(q)$ independently
of earlier and later batches. Dispatch resets only the undispatched backlog;
it does not return vehicles to their home cities or erase goods already in
service. Services of different batches do not compete for a finite fleet or
sorting resource. If costs are instead recorded when physical services occur,
the difference from this prepaid accounting must be $o(T)$ over a horizon
of length $T$ for the same long-run comparison to apply.
\end{assumption}

Let $Q(t)$ be the undispatched backlog, with $Q(0)=0$. Between dispatches,
$\dot Q=\Lambda$. An intervention at a positive backlog dispatches all of
it, incurs $C(Q(t-))$, and sets $Q(t)=0$. No rejection, partial dispatch,
empty dispatch, or additional fleet or age state is permitted. The
information consists of the deterministic arrival rate and past actions.
The horizon $T$ below is elapsed physical time in minutes, distinct from
the integer operating endpoint of Section~\ref{sec:model}.
An admissible policy chooses positive wave sizes $q_k$ with
$\sum_k q_k=\infty$, giving intervention times
$t_k=\sum_{r=1}^k q_r/\Lambda\uparrow\infty$. It therefore has no finite-time
accumulation of interventions. Define
\begin{equation}
 A_\pi(T)=\sum_{t_k\le T}C(q_k)+\int_0^T\omega Q(t)\,dt,
 \qquad
 \inf_\pi\liminf_{T\to\infty}\frac{A_\pi(T)}T.
 \label{eq:analysis-average-control}
\end{equation}
This is an impulse problem because the controller chooses isolated dispatch
times and pays a batch-dependent intervention cost. It is distinct from a
fixed-step rolling-horizon controller.

\begin{theorem}[Static and dynamic common-wave consistency]
\label{thm:analysis-consistency}
Under Assumptions~\ref{ass:analysis-benchmark} and
\ref{ass:analysis-accounting}, the average cost rate of repeatedly dispatching
at a positive threshold $q$ is
\begin{equation}
 J(q)=\frac{\Lambda C(q)}q+\frac{\omega q}2
     =\Lambda G^\circ(q).
 \label{eq:analysis-consistency}
\end{equation}
The set $\mathcal B=\argmin_{q>0}J(q)$ is nonempty and compact. If
$g=\min_{q>0}J(q)$, every admissible policy in
\eqref{eq:analysis-average-control} has limiting inferior average cost at
least $g$, and repeating any threshold in $\mathcal B$ attains $g$.
Thus $\argmin_{q>0}G^\circ(q)=\mathcal B$.
\end{theorem}

The cycle cost is $C(q)+\omega q^2/(2\Lambda)$ and its duration is
$q/\Lambda$. Their ratio gives \eqref{eq:analysis-consistency}.
Appendix~\ref{app:dispatch} proves existence and the lower bound for arbitrary
calendar horizons, including the unfinished final cycle. Thus varying
successive complete-wave sizes cannot improve the best repeated-wave
cost under the stated assumptions; the minimizing wave need not be unique.
For two symmetric directions, both costs and arrivals double, leaving the
threshold set unchanged.

\subsection{Scope and verification summary}

Given $q$, Proposition~\ref{prop:analysis-loads} supplies the optimized
loads and hence $G^\circ(q)$; its global minimizer set characterizes
optimal repeated waves. Appendix~\ref{subsec:qvi} verifies the same
backlog-reset model through a QVI, with a tail-envelope relative value,
the waiting boundary at zero, and a proved transversality condition.
Appendix~\ref{subsec:analysis-limits} also compares unit-Poisson and fluid
costs within the same positive integer full-reset threshold class.
These results retain the benchmark's common-wave and separable-cost
assumptions. They do not prescribe the finite-fleet rolling controller.

\section{Discussion and conclusions}
\label{sec:discussion}

The model connects a local routing approximation to finite-fleet intercity
scheduling. Wave-capacity tables feed a time-expanded MILP, whose vehicle
inventories and freight precedence constraints coordinate collection,
line-haul, distribution, direct service, and independent empty repositioning.
Rolling execution preserves the resulting whole-task commitments and their
inventory consequences across computational windows. A physical terminal
fleet target is a separate operating requirement.

The analytical results explain two different choices. Pooling vehicles on
the same local arc weakly increases modeled capacity, while the conditional
load formulas balance repeated stem travel against shipment lead time.
The collection and distribution loads differ because their delivery-time
exposure differs. In the stationary benchmark, a per-shipment cost minimizer
also determines an average-cost-optimal repeated wave among the stated
successive complete-wave policies. The full calendar-time argument and QVI
verification are retained in Appendix~\ref{app:dispatch}.

These interpretations depend on the declared approximations. Balanced
zoning and stop-equivalent demand require appropriate calibration for an
application; the aggregate direct duration must represent the vehicle's
complete task before it is credited to the destination inventory. The
stationary results additionally require common waves and separable batch
costs, whereas shared fleets, finite sorting, and independently timed
stages can couple batches. Table~\ref{tab:model-scope} keeps these analytical
assumptions distinct from the scheduling model. The implementation reports
partial execution when continuation fails and uses window-level solver
certificates with the scope stated in Appendix~\ref{app:replay}.


\appendix
\section{Capacity derivations and loading proofs}
\label{app:proofs}

\subsection{Local capacity}

Fix a city and suppress its index. Write $f(q)=aq+b\sqrt q$, where $a>0$,
$b\ge0$, and let $D=(j-i)\Delta t$ be the physical duration of a local-service
arc. With $B=(D-2\rho)^+$, the coefficient used in
Section~\ref{sec:capacity} is
\begin{equation}
 K(v)=\min\{vM,f^{-1}(vB)\},\qquad K(0)=0.
 \label{eq:analysis-capacity}
\end{equation}
The inverse of $f$ exists on $[0,\infty)$ and is
\begin{equation}
 f^{-1}(z)=\left(\frac{\sqrt{b^2+4az}-b}{2a}\right)^2.
 \label{eq:analysis-inverse}
\end{equation}
For a positive wave $q$ served by $v>0$ vehicles, $q\le K(v)$ is equivalent to
the load bound $q/v\le M$ and the duration bound
\begin{equation}
 2\rho+\frac{f(q)}v
 =2\rho+\left(a+\frac b{\sqrt q}\right)n\le D,
 \qquad n=\frac qv.
 \label{eq:analysis-local-bridge}
\end{equation}
Thus the per-stop routing term depends on the total wave $q$, whereas the stem
is incurred by every vehicle. This is the zoning mechanism represented by the
BHH-inspired approximation \cite{Beardwood1959}; it is a modeling approximation
for finite tours, rather than an exact routing-length theorem.

\begin{proof}[Proof of Proposition~\ref{prop:analysis-capacity}]
The function $f$ is continuous and strictly increasing on $[0,\infty)$.
Setting $r=\sqrt q$ in $ar^2+br=z$ gives
\eqref{eq:analysis-inverse}. If $B=0$, then $K(v)=0$ for all $v$ and the
claim is immediate. Suppose $B>0$ and write $H=f^{-1}$. For $z>0$,
\[
 \frac{H(z)}z=\frac1{a+b/\sqrt{H(z)}}
\]
is nondecreasing in $z$. Hence
\[
 \frac{K(v)}v=\min\left\{M,
                  B\frac{H(vB)}{vB}\right\}
\]
is nondecreasing in $v>0$. For positive $v_1,v_2$, multiply
$K(v_r)/v_r\le K(v_1+v_2)/(v_1+v_2)$ by $v_r$, for $r=1,2$, and add.
If either argument is zero, superadditivity follows from $K(0)=0$.

For completeness, neither convexity nor strictness follows from this proof.
With $a=b=B=1$ and $M=1/2$,
\[
 K(1)=\frac{3-\sqrt5}2,\qquad K(2)=1,\qquad K(3)=\frac32.
\]
Thus $K(3)-2K(2)+K(1)<0$, which rules out convexity even on the integer
domain. On the capped region $K(v)=vM$, so equality can hold in
\eqref{eq:analysis-superadditive}. Zero budget and $b=0$ also remove strict
routing economies.
\end{proof}

\subsection{Conditional vehicle loads}

\begin{proof}[Proof of Proposition~\ref{prop:analysis-loads}]
For fixed $q>0$, terms involving the loads in
\eqref{eq:analysis-generalized-cost} are
\[
 \frac{2c_H\rho}{n_1}+\omega s(q)n_1
 +\frac{2c_H\rho}{n_2}+\frac{\omega s(q)n_2}2.
\]
Each summand is strictly convex in its positive load and diverges as that
load decreases to zero. Its unconstrained minimizer is the corresponding
square root in \eqref{eq:analysis-n1}--\eqref{eq:analysis-n2}. Clipping at
the common upper bound $U=\min(M,q)$ gives the unique constrained minimizers.
Writing $r=\sqrt{2c_H\rho/(\omega s(q))}$, they are $\min(U,r)$ and
$\min(U,\sqrt2r)$. Their ordering and ratio bound follow directly.
\end{proof}


\section{Calendar-time optimality and dispatch verification}
\label{app:dispatch}

\subsection{Existence and calendar-horizon optimality}

We first record regularity and growth properties needed in both verification
arguments. The clipped loads are continuous and locally Lipschitz for
$q>0$. Thus $C$ and $J$ are continuous and locally Lipschitz there, including
across load and highway-capacity switches. Since $n_i\le q$, the stem terms
in the batch service cost imply
\begin{equation}
 C(q)\ge4c_H\rho>0,
 \qquad
 J(q)\ge\frac{4c_H\rho\Lambda}q+\frac{\omega q}2.
 \label{eq:proof-growth-lower}
\end{equation}
In particular $J$ diverges at zero and at infinity. It attains its minimum
$g$ on a compact subset of $(0,\infty)$, and its minimizer set
$\mathcal B$ is nonempty and compact. This argument does not require
convexity of the reduced objective or uniqueness of its minimizer.

There is also a constant $K_0<\infty$ such that
\begin{equation}
 0\le C(q)\le K_0(1+q),\qquad q>0.
 \label{eq:proof-growth-upper}
\end{equation}
For $q\ge1$, evaluate the unoptimized service cost at the fixed feasible
loads $n_1=n_2=\min(M,1)$. Then $s(q)\le a+b$, and the per-shipment
service cost is bounded. For $0<q\le1$, the feasible choice
$n_1=n_2=\min(M,q)$ gives a bounded batch cost after multiplication by
$q$: the stems are bounded, $qs(q)=aq+b\sqrt q$ is bounded, and
$q/\min(q,M_A)$ is bounded. Optimizing the loads can only reduce this
upper bound. In fact, the clipped optimizers equal $q$ for sufficiently
small $q$, and
\begin{equation}
 \lim_{q\downarrow0}C(q)=4c_H\rho+c_A\tau_{\mathrm{phys}}.
 \label{eq:proof-zero-limit}
\end{equation}

For a policy in Section~\ref{subsec:dynamic-wave}, let $r$ be the elapsed
time since its last dispatch by $T$. The calendar-time bound is
\begin{equation}
 A_\pi(T)-gT
 =\sum_{k:t_k\le T}\frac{q_k}{\Lambda}[J(q_k)-g]
     +\frac{\omega\Lambda r^2}2-gr
 \ge-\frac{g^2}{2\omega\Lambda}.
 \label{eq:analysis-calendar-bound}
\end{equation}

\begin{proof}[Proof of Theorem~\ref{thm:analysis-consistency}]
Starting immediately after a dispatch, the backlog grows from zero to $q$
in $q/\Lambda$ minutes. Its holding cost over that cycle is
$\omega q^2/(2\Lambda)$. Adding $C(q)$ and dividing by the cycle length
proves \eqref{eq:analysis-consistency}. Existence and compactness follow
from \eqref{eq:proof-growth-lower}.

Consider any admissible policy and a calendar time $T$. Admissibility
ensures that the number $N=N(T)$ of dispatches by $T$ is finite. Write
$r=T-t_N$, taking $t_0=0$ if $N=0$. Each complete cycle contributes
$(q_k/\Lambda)J(q_k)$; the unfinished cycle contributes
$\omega\Lambda r^2/2$. Subtracting $gT$ gives
\eqref{eq:analysis-calendar-bound}, because all completed-cycle excesses
are nonnegative and
\[
 \frac{\omega\Lambda r^2}2-gr
 =\frac{\omega\Lambda}2
        \left(r-\frac g{\omega\Lambda}\right)^2
       -\frac{g^2}{2\omega\Lambda}.
\]
This uniform lower bound proves the limiting inferior claim. If a threshold
$b\in\mathcal B$ is repeated, all complete-cycle excesses vanish and
$0\le r<b/\Lambda$. The residual term divided by $T$ tends to zero, so
the complete calendar-time average converges to $g$. Multiplication by
the positive constant $\Lambda$ preserves the static minimizer set.
\end{proof}

\subsection{A verified average-cost QVI on the reduced state space}
\label{subsec:qvi}

Normalize the relative value by $V(0)=0$. For $q>0$, the intervention
operator is $\mathcal MV(q)=C(q)+V(0)$. The average-cost
quasi-variational inequality is
\begin{equation}
 \min\left\{\omega q-g+\Lambda V'(q),\ C(q)-V(q)\right\}=0.
 \label{eq:analysis-qvi}
\end{equation}
This is complementarity notation: the waiting cost-rate residual and
the intervention cost residual are both nonnegative, with at least one
zero. At $q=0$
dispatch is unavailable, and the waiting equation supplies the boundary
condition $\Lambda V'(0+)=g$. The link with impulse control is the
intervention operator; QVI terminology alone does not supply verification
or an optimal policy \cite{LionsPerthame1986,Belak2017}.

In the construction below, $F$ denotes an excess-cost function, distinct
from the city routing inverse $F_k$ in Section~\ref{sec:capacity}.

\begin{theorem}[Global relative-value construction and verification]
\label{thm:analysis-qvi}
Let $g$ be as in Theorem~\ref{thm:analysis-consistency} and define
\begin{align}
 F(q)&=C(q)+\frac{\omega q^2}{2\Lambda}-\frac{gq}{\Lambda},
 &m(q)&=\inf_{z\ge q}F(z),\qquad q>0,
 \label{eq:analysis-tail-envelope}\\
 V(q)&=\frac{gq-\omega q^2/2}{\Lambda}+m(q),
 &V(0)&=0.
 \label{eq:analysis-value}
\end{align}
The tail infima are attained. The function $V$ is locally absolutely
continuous, satisfies \eqref{eq:analysis-qvi} almost everywhere on
$(0,\infty)$, and has $V'(0+)=g/\Lambda$. Along every admissible policy,
\begin{equation}
 A_\pi(T)-gT\ge V(Q(0))-V(Q(T)).
 \label{eq:analysis-verification-bound}
\end{equation}
Moreover $V(q)=O(1+q)$ in absolute value. On every sequence $T_j\uparrow\infty$
with bounded $A_\pi(T_j)/T_j$, the transversality term
$V(Q(T_j))/T_j$ tends to zero. This verifies the same average-cost lower
bound and the optimal repeated thresholds.
\end{theorem}

The tail envelope is needed for a statement on all initial backlog states.
The cycle minimizer alone does not show that immediately dispatching at
every state above its smallest value is an optimal first action. Equivalently,
the constructed relative value is the first-dispatch excess-cost envelope
\begin{equation}
 V(q)=\inf_{z\ge q}\left\{C(z)
       +\frac{\omega(z^2-q^2)}{2\Lambda}
       -\frac{g(z-q)}{\Lambda}\right\}.
 \label{eq:analysis-first-dispatch}
\end{equation}
If $b_*:=\min\mathcal B$, then $m(q)=0$ on $[0,b_*]$. At a differentiable
interior optimum $b_*$, value matching and the stationary first-order
condition give $V(b_*)=C(b_*)$ and
$C'(b_*)=(g-\omega b_*)/\Lambda$. At a load or capacity kink these
identities must be interpreted using the available one-sided derivatives;
classical smooth fit is not imposed. No uniqueness of a QVI solution or
general viscosity comparison theorem is asserted here.

\begin{proof}[Proof of Theorem~\ref{thm:analysis-qvi}]
The definition of $g$ gives
\[
 F(q)=\frac q\Lambda[J(q)-g]\ge0,
 \qquad \{q>0:F(q)=0\}=\mathcal B.
\]
The function $F$ is locally Lipschitz on $(0,\infty)$, has the finite
right limit in \eqref{eq:proof-zero-limit}, and tends to infinity as
$q\to\infty$. Its tail infima are therefore attained. It is not
necessary, and generally not correct, to assert that $F$ is Lipschitz
at zero: BHH terms may contribute $\sqrt q$ to $C(q)$.

The tail minimum $m$ is nondecreasing. It is also locally Lipschitz on
$(0,\infty)$. To see this, take a compact interval $[r,R]$ with $r>0$.
Coercivity of $F$ allows all tail minimizers for arguments in $[r,R]$
to be chosen in a fixed compact interval $[r,Z]$. Let $L_F$ be a
Lipschitz constant of $F$ there. For $r\le x<y\le R$, if a minimizer
for $m(x)$ lies at or above $y$, then $m(x)=m(y)$. Otherwise choose
one $z_x\in[x,y)$; since $m(y)\le F(y)$,
\[
 0\le m(y)-m(x)\le F(y)-F(z_x)\le L_F(y-x).
\]
At every point with $F(q)>m(q)$, a tail minimizer lies strictly to its
right. Continuity then makes $m$ constant on a neighborhood of $q$:
a small change in the lower tail boundary neither removes that minimizer
nor admits a smaller value. Hence $m'=0$ almost everywhere on the strict
noncontact region.

Let $b_* =\min\mathcal B>0$. Nonnegativity of $F$ and $F(b_*)=0$ imply
$m(q)=0$ for $0<q\le b_*$. Set $m(0)=0$. The resulting $V$ is a smooth
polynomial near zero, is locally absolutely continuous elsewhere, and
has $V(0)=0$ and $V'(0+)=g/\Lambda$. Almost everywhere on $(0,\infty)$,
\[
 C(q)-V(q)=F(q)-m(q)\ge0,
 \qquad
 \omega q-g+\Lambda V'(q)=\Lambda m'(q)\ge0.
\]
If the first quantity is positive, local constancy of $m$ makes the
second zero. If it is zero, complementarity already holds. This proves
the QVI in the stated almost-everywhere sense. Algebraically expanding
the tail-infimum definition also gives
\eqref{eq:analysis-first-dispatch}.

Integrate the waiting inequality along each uninterrupted drift segment.
At an intervention from backlog $q$ to zero, use
$C(q)\ge V(q)-V(0)$. Summing over the finitely many interventions
before $T$ telescopes all intermediate values of $V$ and proves
\eqref{eq:analysis-verification-bound}. This argument is valid for
locally absolutely continuous $V$, since backlog grows linearly on each
drift segment and thus spends zero time in the exceptional set of
nondifferentiability.

To verify transversality rather than presume it, first use immediate
dispatch in \eqref{eq:analysis-first-dispatch} to obtain $V(q)\le C(q)$.
The same formula and $C\ge0$ imply
\[
 V(q)\ge \inf_{z\ge q}
      \int_q^z\frac{\omega u-g}{\Lambda}\,du
      \ge-\frac{g^2}{2\omega\Lambda}.
\]
Together with \eqref{eq:proof-growth-upper}, this gives
$|V(q)|\le K_1+K_2q$ for finite constants. At any calendar time, the
unfinished cycle alone has incurred holding cost
$\omega Q(T)^2/(2\Lambda)$. All other costs are nonnegative. Therefore,
on any sequence with bounded $A_\pi(T_j)/T_j$,
\[
 Q(T_j)=O(\sqrt{T_j}),\qquad
 \frac{|V(Q(T_j))|}{T_j}\longrightarrow0.
\]
If the limiting inferior average cost is finite, choose such a sequence
realizing it and apply \eqref{eq:analysis-verification-bound}; if it is
infinite, the lower bound is immediate. For a repeated optimal threshold
the backlog is bounded, all relevant inequalities are equalities along
the regenerative path, and the same bound verifies its average cost $g$.
\end{proof}

The construction also describes optimal first-dispatch excess costs from
arbitrary positive backlogs: wait until a point attaining the tail minimum
in \eqref{eq:analysis-first-dispatch}, dispatch, and then repeat any
threshold in $\mathcal B$. For states above $b_*$, immediate dispatch
requires an additional shape property of $F$; for example, it follows if
$F$ is nondecreasing on $[b_*,\infty)$. Merely minimizing $J$ at $b_*$
does not establish that property. At a differentiable interior minimizer,
$F'(b_*)=0$ gives the derivative statement above. At a
nondifferentiable local minimum the valid one-sided statement is
$F'_-(b_*)\le0\le F'_+(b_*)$, when these derivatives exist.

\subsection{Unit Poisson arrivals within the integer-threshold class}
\label{subsec:analysis-limits}

There is one useful exact stochastic comparison. Let arrivals instead be a
unit-jump Poisson process of rate $\Lambda$, with its natural filtration,
and restrict decisions to dispatching all shipments immediately when the
backlog reaches a fixed positive integer $B$. Retain the same $C(B)$ and
linear holding cost. Renewal accounting gives the following result.

\begin{proposition}[A common integer-threshold comparison]
\label{prop:analysis-poisson}
For these Poisson threshold policies,
\begin{equation}
 J_P(B)=\frac{\Lambda C(B)}B+\frac{\omega(B-1)}2
       =J(B)-\frac\omega2,\qquad B\in\{1,2,\ldots\}.
 \label{eq:analysis-poisson}
\end{equation}
Consequently the Poisson and deterministic-fluid problems have the same
minimizing thresholds when both are restricted to this same integer set.
Their optimal cost rates differ by $\omega/2$.
\end{proposition}

This restricted result is not an optimality theorem over all adapted
stochastic intervention policies. The Poisson waiting generator is
$\Lambda[V(n+1)-V(n)]$, rather than $\Lambda V'(q)$; a general stochastic
verification requires integrability and expected transversality conditions.
The continuous-wave and integer-wave domains also need not have the same
optimizer. For $B>1$, the time required to collect the remaining Poisson
arrivals has no deterministic upper bound, whereas $B=1$ dispatches at
arrival and creates no consolidation wait.

\begin{proof}[Proof of Proposition~\ref{prop:analysis-poisson}]
For a fixed integer $B\ge1$, the dispatch cycle contains exactly $B$
independent exponential interarrival times of mean $1/\Lambda$. Its
mean duration is $B/\Lambda$. The backlog spends expected time
$1/\Lambda$ in each state $0,1,\ldots,B-1$, so expected holding cost
per cycle is
\[
 \frac\omega\Lambda\sum_{n=0}^{B-1}n
       =\frac{\omega B(B-1)}{2\Lambda}.
\]
The service cost per cycle is $C(B)$. The renewal-reward ratio is
therefore \eqref{eq:analysis-poisson}. The difference from the fluid
criterion is the same constant for every admissible integer threshold,
which preserves its minimizer set. Both integer minima exist, since
the holding term tends to infinity with $B$. This calculation proves
the claim within the stated threshold class and makes no comparison
with other adapted intervention rules.
\end{proof}

The restriction to a common action domain is substantive. For example,
in a separate setup-cost illustration with $C(q)=9/8$ and
$\Lambda=\omega=1$, the continuous-fluid optimum is $q=3/2$ with cost
rate $3/2$. The same fluid criterion restricted to positive integers
has its unique minimum at $B=2$, with cost rate $25/16$. The Poisson
threshold criterion also chooses $B=2$ and has cost rate $17/16$.
Thus equality of mean arrival rates alone implies neither equality
with the continuous optimizer nor equality of cost rates. This example
illustrates the action-domain issue; its constant setup cost is not
the BHH service cost in \eqref{eq:analysis-dispatch-cost}.


\section{Candidate filtering and independent schedule replay}
\label{app:implementation}

\subsection{Time-support filtering of cargo variables}
\label{app:filter}

The sparse variant removes cargo variables, while keeping the vehicle and
wave-selector mother domains. It therefore preserves automated repositioning,
independent empty transfers, and zero-cargo vehicle movements. The filter
does not infer future city availability from historical idle stock alone:
a vehicle that is absent now may arrive before a future task.

For a revealed order $\ell\in\mathcal L^d\cap\mathcal L(t)$, form provisional new
collection candidates $P_\ell$ obeying \eqref{alg:origin-cut},
$i\geq s_\ell$, and positive local capacity at the upper fleet size $N$.
Form provisional new distribution candidates $R_\ell$ obeying
\eqref{alg:downstream-cut}, $j\leq e_\ell$, and positive local capacity.
Let $J^P_\ell$ be the end times of committed positive collection cargo and
$J^A_\ell$ the arrival times of committed positive automated cargo. These
history sets seed downstream possibilities even when there is no admissible
new collection tour. Define
\begin{equation}
 \underline t_\ell^P=\min\bigl(\{j:(i,j)\in P_\ell\}\cup J^P_\ell\bigr),
 \qquad
 \overline t_\ell^R=\max\{i:(i,j)\in R_\ell\},
 \label{alg:time-bounds}
\end{equation}
with the empty-set values $+\infty$ and $-\infty$, respectively. If
$A>0$, retain provisional automated cargo arcs
\begin{equation}
 Q_\ell=\{(i,j)\in\mathcal A_A:
 t\leq i<c_t,\ j\leq c_t,\quad
 \underline t_\ell^P\leq i,\ j\leq\overline t_\ell^R\};
 \label{alg:auto-support}
\end{equation}
if $A=0$, take $Q_\ell=\varnothing$. A new collection cargo variable is
retained only if its end is no later than some start in $Q_\ell$.
A new distribution cargo variable is retained only if its start is at least
\begin{equation}
 \min\bigl(\{j:(i,j)\in Q_\ell\}\cup J^A_\ell\bigr).
 \label{alg:delivery-support}
\end{equation}
Direct cargo candidates must obey $i\geq\max\{t,s_\ell\}$,
$i<b_t$, $j\leq\min\{c_t,e_\ell\}$, and
$\widetilde K^d_{ij}(N)>0$. If $N=0$, all new human-driven cargo
candidates are empty. All previously committed nonzero cargo keys are
retained regardless of these tests so that history can be fixed exactly.

These tests use time support and optimistic capacity; they do not solve
the joint fleet-allocation problem. Retaining a candidate is therefore no
claim that the candidate is feasible together with other decisions. In
the MILP, omitted cargo variables are zero, and every remaining cargo
variable is still subject to the original capacity, vehicle, precedence,
and demand constraints. The procedure records the number of cargo
variables before and after filtering for each rebuilt window. We do not
assert a general preservation theorem or a computational speed advantage
from this filter; neither follows from the construction alone.

\subsection{Independent replay and certificate scope}
\label{app:replay}

A returned schedule is checked independently of the optimizer's capacity
tables. For example, a positive local wave load $q$ and count $v$ must
satisfy
\begin{equation}
 q\leq vM,\qquad
 f_k(q)\leq v[(j-i)\Delta t-2\rho_k]^+,
 \label{alg:local-replay}
\end{equation}
and positive direct cargo must satisfy $q\leq vM$ and
\eqref{mod:direct-duration} within its assigned duration. Replay also
checks domains, selector links, shared automated capacity, every city
inventory, order windows, cumulative precedence, and demand balance.
Costs are recomputed from initiated whole tasks once, rather than by
summing overlapping window objective values.

For a completed offline solve, a reported MILP gap refers only to the
specified finite-horizon model and its input information. Window solver
bounds and gaps refer to their individual restricted subproblems; the
last window's bound is not a global lower bound for rolling execution.
For partial execution, replay can certify past commitments and accounting,
including hub or in-transit goods, but defers final stage balance and future
completion reachability. A partial schedule consequently has no certified
complete-horizon objective. These distinctions are also needed when
comparing this finite-fleet dispatch procedure with the reduced
steady-state and QVI models.



\bibliographystyle{plainnat}
\bibliography{references}
\end{document}
